\section{Introduction}
Fraud rings are groups of accounts that act together: they share devices,
payment cards or network addresses, and they move money among themselves in
cycles or fan-in/fan-out patterns. Property graphs make these relations
explicit, and graph pattern queries are a common way to surface them
\cite{pourhabibi2020fraud,akoglu2015graph}. Two problems appear as soon as such
queries run on a live stream of events.

\textbf{The hub dilemma.} Legitimate shared infrastructure, such as an office
gateway, a carrier-grade NAT address or a store kiosk, links hundreds of
unrelated accounts. A query that treats every shared attribute as evidence
floods analysts with false rings. The usual remedy is a hard degree cap:
attributes shared by more than $c$ accounts are ignored. The cap removes the
noise but also opens an evasion path, because a ring that shares a busy
attribute becomes invisible to it.

\textbf{Latency on a stream.} Accounts, links and transfers arrive
continuously. Re-running whole-graph queries on a timer delays detection, and
the cost of each pass grows with the graph rather than with the change.

We make four contributions.
\begin{itemize}
\item \textbf{C1, hub-aware weighting.} We replace the degree cap with a
continuous weight: each shared attribute contributes $1/\ln(1+\deg)$ with
exponential time decay, and transfers inside the pair add weight. Busy hubs
count for little on their own, but a group that shares a hub and also trades
among itself still accumulates weight (Section~\ref{sec:method}).
\item \textbf{C2, incremental detection.} A micro-batch engine re-evaluates only
the $k$-hop region touched by each batch, with fan-out limits that keep that
region bounded when it reaches a hub.
\item \textbf{C3, calibrated, explained alerts.} Candidate rings are scored by a
calibrated classifier, and each alert carries feature contributions, an
evidence subgraph and counterfactuals.
\item \textbf{C4, evaluation under evasion.} A deterministic simulator with
ring-level ground truth and five evasion strategies, including an attacker that
adapts to the detector's flags, lets us measure each component on identical
streams (paired design) with confidence intervals.
\end{itemize}

\section{Related Work}
\label{sec:related}
\textbf{Graph-based fraud detection.} Surveys of graph anomaly detection
\cite{akoglu2015graph} and of graph-based fraud detection
\cite{pourhabibi2020fraud} cover dense-subgraph, community and pattern-based
methods. Dense-subgraph methods such as FRAUDAR \cite{hooi2016fraudar} are
explicitly designed to resist camouflage edges; they operate on static
bipartite graphs, while we target a stream with ring-level alerts.
\TODOcite{dense subgraph detection on streams (e.g. incremental densest subgraph)}

\textbf{Community detection.} Louvain modularity optimisation
\cite{blondel2008louvain} is a common unsupervised baseline for finding groups
of colluding accounts; we use it as a baseline.

\textbf{Streaming and incremental graph analytics.}
\TODOcite{incremental graph pattern matching / continuous subgraph queries}
\TODOcite{Neo4j GDS or streaming graph database systems}

\textbf{Graph neural networks for AML.} GNNs have been applied to transaction
graphs such as Elliptic \cite{weber2019elliptic} and AMLworld
\cite{altman2023amlworld}. They classify transactions or accounts rather than
emit groups, and they are not evaluated here (Section~\ref{sec:threats}).
\TODOcite{GNN fraud ring / camouflage-resistant GNN work}

\textbf{Hub handling.} Down-weighting high-degree nodes by
$1/\log(\deg)$ follows the Adamic--Adar idea from link prediction.
\TODOcite{Adamic and Adar 2003, Friends and neighbors on the Web}

\section{Method}
\label{sec:method}
Fig.~\ref{fig:arch} shows the system. A seeded simulator emits observable
events (account joins and leaves, attribute links, transfers); ground truth
goes only to a ledger the detector never sees.

\begin{figure}[t]
\centering
\includegraphics[width=\columnwidth]{\figpath/architecture.pdf}
\caption{Architecture. Solid arrows are the observable data path; dashed
arrows carry labels (to evaluation only) and, in closed-loop runs, the
detector's flags back to the adaptive attacker.}
\label{fig:arch}
\end{figure}

\subsection{Graph model}
Users link to devices, cards and IP addresses
(\texttt{USED\_DEVICE}, \texttt{HAS\_CARD}, \texttt{LOGGED\_IN\_FROM}) and send
transfers to each other. The same contract is implemented over Neo4j (Cypher)
and over an in-memory store with identical semantics; all experiments below use
the in-memory store.

\subsection{Hub-aware weighting (C1)}
For two users $u,v$ the edge weight is
\begin{equation}
W(u,v)=\sum_{a\in A(u)\cap A(v)} h(\deg a)\,\delta(t_a)
+\min\Big(\beta\sum_{x\in T(u,v)}\delta(t_x),\,\beta_{\max}\Big),
\label{eq:w}
\end{equation}
where $h(d)=1/\ln(1+d)$, $\delta(t)=2^{-t/\tau}$ is exponential decay with
half-life $\tau=300$\,s, $t_a$ is the age of the older of the two links to
attribute $a$, and $T(u,v)$ are the transfers between $u$ and $v$ in the
window, with $\beta=0.15$ and $\beta_{\max}=0.6$. A device shared by three
accounts weighs $0.72$; an office address shared by fifty weighs $0.25$.
Candidate rings are the connected components of the graph of pairs with
$W\ge\theta_e$ (and a non-zero attribute term), peeled to cores of edge density
at least $0.75$ and size 3--25. The edge threshold $\theta_e$ is selected on
validation seeds (Section~\ref{sec:protocol}).

\subsection{Incremental detection (C2)}
Events are grouped into micro-batches of 50. For each batch the engine takes
the users it touched, expands them by $k=2$ hops through shared attributes and
transfers, and builds a view of that region over a 900\,s window, together with
the \emph{global} degree of every attribute in it, so $h(\deg a)$ is the same as
in a full pass. The first hop may cross attributes with up to 250 users (so the
co-members of a hub are re-examined); later hops cross only attributes with at
most 20 users, which stops the region from spreading through hubs. Incremental
detectors (hub-aware shared attribute, temporal cycle, smurfing) run on the
view; whole-graph detectors (the static Cypher rules and community detection)
run on a timer.

\subsection{Ensemble, tracking and scoring (C3)}
Candidates from different detectors are merged when their Jaccard overlap is at
least $0.5$, or when one contains $80\%$ of the other and they differ in size by
at most a factor of two. A tracker matches candidates to existing rings across
batches; a ring is \emph{confirmed} after two consecutive hits and dissolved
after three misses. Confirmed rings are scored by a logistic-regression model
over 19 ring features (size, density, which attribute kinds are shared, the
hub-adjusted weight, cycle presence and length, burstiness, account ages,
internal transfers per member, amount statistics, growth, hub exposure and
the number of agreeing detectors), calibrated with isotonic regression
\cite{zadrozny2002isotonic}. A ring becomes an alert when its calibrated risk
reaches the alert threshold $\theta_a$. Each alert carries the per-feature
contributions to the score, the subgraph of shared attributes and transfers
that formed it, and counterfactuals (for the leading features, the value at which
the risk would fall below the threshold, and for hub-aware rings the degree the
shared attribute would need before the rule stops firing). The model is trained on
seeds disjoint from all evaluation seeds.

\section{Experimental Setup}
\label{sec:setup}
\subsection{Simulator and scenarios}
The simulator is a deterministic discrete-event generator: a scenario file and
a seed fix the whole stream, which we verify by a digest of the event sequence
recorded with every result. Benign users join and leave, link to devices, cards
and IPs drawn from heavy-tailed popularity, and transact; benign hubs (offices,
carrier NAT, kiosks, family devices) are generated with configured sizes. Ring
types are shared device, card or IP, mixed attribute, money cycle, layered
chain and smurfing fan-in/fan-out. Evasion strategies are hub hiding (the ring
shares a busy benign hub), low-and-slow activity, identity rotation, decoy edges
to benign users, and an adaptive attacker that shifts its strategy mix away
from strategies that were flagged.
Table~\ref{tab:scenario-def} lists the scenarios.

\begin{table}[t]
\centering
\caption{Scenarios (from \texttt{configs/scenarios/*.yaml}).}
\label{tab:scenario-def}
\small
\setlength{\tabcolsep}{3pt}
\begin{tabular}{lrrrl}
\toprule
Scenario & Users & Dur.\,(s) & Ev./s & Rings and evasion \\
\midrule
baseline & 30 & 60 & 10 & 4 planted, none \\
high\_noise & 150 & 60 & 50 & 4 planted + random, decoys \\
hub\_heavy & 100 & 60 & 25 & 4 planted, hub hiding \\
adversarial & 120 & 90 & 30 & 5 planted + random, all five \\
drift & 200 & 90 & 40 & 2 planted + random, hub hiding \\
demo\_story & 120 & 90 & 20 & 3 planted, hub hiding \\
scale\_sweep & $10^3$+ & 120 & 100 & latency benchmark (E4) \\
\bottomrule
\end{tabular}
\end{table}

\subsection{Protocol}
\label{sec:protocol}
Seeds are split three ways: scorer training (101--110, test 201--205),
threshold selection (validation, 301--303) and evaluation (test, 1001 and up).
On the validation seeds we evaluated a grid of alert thresholds
$\theta_a\in\{0.2,0.3,0.5,0.7,0.8,0.9\}$ and edge thresholds
$\theta_e\in\{0.2,0.3,0.4,0.5,0.6\}$ on the five accuracy scenarios; the grid was
widened whenever the best value sat on its edge. The selected values are
$\theta_a=\TunedAlert$ and $\theta_e=\TunedEdge$ (mean validation F1
$\TunedValF$). We report the proposed method both with these values and with
the a-priori values $\theta_a=\theta_e=0.5$.

Every method in an experiment processes the \emph{identical} event stream for a
given scenario and seed (the attacker does not receive feedback in these runs),
so differences between methods are paired. We report the mean and 95\%
$t$-interval over runs, and test differences with the Wilcoxon signed-rank test
on paired runs, with Cohen's $d_z$ as effect size.

\subsection{Metrics}
A flagged ring \emph{matches} a true ring when their member sets overlap with
Jaccard at least $0.5$. Ring precision is the fraction of flagged rings that
match a true ring; ring recall is the fraction of true rings matched at some
point while active. A run in which nothing is flagged has undefined precision;
we count its F1 as 0. Time to detect (TTD) is the time from a ring's
activation to its first matching flag. Detection latency is the wall-clock time
per micro-batch for the dirty-region expansion, the view and the incremental
detectors.

\subsection{Baselines}
\begin{itemize}
\item \textbf{Static Cypher, cap 100 / cap 10.} The original rules: shared
device (Q1), shared card (Q2), 3-cycles of transfers (Q3), and shared IP with a
degree cap (Q4a/Q4b), run on the whole graph every 5\,s of simulated time.
\item \textbf{Louvain.} Louvain communities \cite{blondel2008louvain} of the
user graph projected through shared attributes and transfers, kept when dense
enough.
\item \textbf{Connected components.} Weakly connected components of the same
projection, no density filter.
\end{itemize}
All baselines use the same two-hit confirmation as the proposed method.

\subsection{Environment}
All runs used one laptop (Intel Core i3-1115G4, 2 cores / 4 threads, 7.7\,GB
RAM, Windows 11), Python 3.11.9, NumPy 2.2.6, scikit-learn 1.7.2,
NetworkX 3.5, with experiments run one at a time.

\section{Results}
\label{sec:results}

\subsection{RQ1: accuracy against baselines}
% generated from results/e1_baselines/results.csv (sha256 7725d8fa88a896e5, 300 rows)
\begin{table*}[t]
\centering
\caption{Ring-level detection on 5 scenarios $\times$ 10 test seeds (mean $\pm$ 95\% CI over runs;
a run with no flags has undefined precision and counts as F1 $=0$). TTD: median time from ring activation to
first flag. Latency: median per-batch detection time on the in-memory store; $^\dagger$full-graph
methods run on a 5\,s timer, not per batch, so their entry is the median time of one full-graph pass. $p$: paired Wilcoxon signed-rank
test of F1 against the proposed method on identical streams.}
\label{tab:baselines}
\begin{tabular}{lcccccc}
\toprule
Method & Precision & Recall & F1 & TTD (s) & Latency (ms) & $p$ \\
\midrule
Static Cypher (cap 100) & 0.12 $\pm$ 0.03 & 0.61 $\pm$ 0.08 & 0.19 $\pm$ 0.05 & 9.5 & 6.0$^\dagger$ & $<$0.001 \\
Static Cypher (cap 10) & 0.14 $\pm$ 0.04 & 0.61 $\pm$ 0.08 & 0.21 $\pm$ 0.05 & 9.5 & 5.9$^\dagger$ & $<$0.001 \\
Louvain communities & 0.39 $\pm$ 0.09 & 0.33 $\pm$ 0.08 & 0.29 $\pm$ 0.06 & 18.8 & 38.4$^\dagger$ & $<$0.001 \\
Connected components & -- & 0.00 $\pm$ 0.00 & 0.00 $\pm$ 0.00 & -- & 11.1$^\dagger$ & $<$0.001 \\
Hub-aware incremental (proposed) & 0.76 $\pm$ 0.06 & 0.74 $\pm$ 0.06 & \textbf{0.73 $\pm$ 0.05} & 9.3 & 26.6 & -- \\
Proposed, a-priori thresholds & 0.73 $\pm$ 0.06 & 0.69 $\pm$ 0.06 & 0.68 $\pm$ 0.05 & 9.0 & 22.7 & 0.034 \\
\bottomrule
\end{tabular}
\end{table*}

% generated from results/e1_baselines/results.csv (sha256 7725d8fa88a896e5, 300 rows)
\begin{table*}[t]
\centering
\caption{Ring F1 per scenario (mean $\pm$ 95\% CI over 10 test seeds).}
\label{tab:scenarios}
\begin{tabular}{lccccc}
\toprule
Method & baseline & high\_noise & hub\_heavy & adversarial & drift \\
\midrule
Static Cypher (cap 100) & 0.48 $\pm$ 0.07 & 0.06 $\pm$ 0.02 & 0.24 $\pm$ 0.04 & 0.14 $\pm$ 0.04 & 0.05 $\pm$ 0.02 \\
Static Cypher (cap 10) & 0.54 $\pm$ 0.08 & 0.07 $\pm$ 0.02 & 0.27 $\pm$ 0.05 & 0.14 $\pm$ 0.04 & 0.05 $\pm$ 0.02 \\
Louvain communities & 0.44 $\pm$ 0.09 & 0.26 $\pm$ 0.19 & 0.24 $\pm$ 0.09 & 0.14 $\pm$ 0.09 & 0.38 $\pm$ 0.16 \\
Connected components & 0.00 $\pm$ 0.00 & 0.00 $\pm$ 0.00 & 0.00 $\pm$ 0.00 & 0.00 $\pm$ 0.00 & 0.00 $\pm$ 0.00 \\
Hub-aware incremental (proposed) & 0.99 $\pm$ 0.03 & 0.62 $\pm$ 0.08 & 0.79 $\pm$ 0.08 & 0.59 $\pm$ 0.14 & 0.67 $\pm$ 0.04 \\
Proposed, a-priori thresholds & 0.93 $\pm$ 0.07 & 0.66 $\pm$ 0.05 & 0.66 $\pm$ 0.12 & 0.53 $\pm$ 0.11 & 0.62 $\pm$ 0.09 \\
\bottomrule
\end{tabular}
\end{table*}


Table~\ref{tab:baselines} pools \ERuns{} runs (\EScenarios{} scenarios
$\times$ \ESeeds{} test seeds $\times$ six methods). The proposed method reaches
ring F1 \EPropF{} ($\pm$\EPropFCI{}), against \ECapTenF{} for the static rules
with a degree cap of 10, \ECapHundredF{} with a cap of 100, \ELouvainF{} for
Louvain and \EPlainCCF{} for connected components; each difference is
significant on the paired streams ($p<0.001$). The static rules reach recall
\ECapTenR{} but precision only \ECapTenP{}: the shared-attribute rules and the
3-cycle rule fire on many benign groups, and two-hit confirmation does not
remove them. Connected components never produce a group within the ring size
limits, because benign hubs join almost all users into one component; Louvain
splits that component, but its communities rarely coincide with a ring. The
thresholds selected on validation seeds improve on the a-priori ones (F1
\EPropF{} vs.\ \EPropDefaultF{}, $p=0.03$).

Per scenario (Table~\ref{tab:scenarios}), the gap to the static rules is
largest on \texttt{drift} and \texttt{high\_noise}, where many benign groups
satisfy the rules and their precision collapses, and smallest on the clean
\texttt{baseline} scenario and on \texttt{adversarial}, which is the hardest
scenario for every method.

\textbf{Degree cap vs.\ hub weighting.} Fig.~\ref{fig:frontier} isolates the
mechanism this paper replaces. Only shared-attribute and mixed rings count as
ground truth here (flags that match only other ring types are ignored), the
static rules run without the cycle rule, and the hub-aware detector runs alone.
Without the scorer, no degree cap reaches a recall above \FrCapMaxR{}: rings
that share a busy attribute are invisible to every cap. Hub weighting reaches
recall \FrHubMaxR{}, but on its own at precision \FrHubPAtMaxR{}. With the
calibrated scorer on top, hub weighting reaches recall \FrHubSMaxR{} (precision
\FrHubSPAtMaxR{}) and its best F1 is \FrHubSBestF{} (alert threshold
\FrHubSBestFAt{}); the cap with the scorer stays below recall \FrCapSMaxR{} and
its best F1 is \FrCapSBestF{}. At low recall the cap with the scorer is the
more precise of the two (\FrCapSBestFP{} at recall \FrCapSBestFR{}). Hub
weighting therefore does not dominate the cap everywhere; it extends the
frontier to recall levels the cap cannot reach, and the scorer is what makes
those levels usable.

\begin{figure}[t]
\centering
\includegraphics[width=\columnwidth]{\figpath/fig1_pr_frontier.pdf}
\caption{Precision and recall on shared-attribute rings (micro-averaged over
4 scenarios $\times$ 5 test seeds). Solid: detector alone, swept over the
degree cap or the edge threshold $\theta_e$. Dashed: with the calibrated scorer,
swept over the alert threshold $\theta_a$ (cap 10; $\theta_e=\TunedEdge$).}
\label{fig:frontier}
\end{figure}

\subsection{RQ2: latency and scale}
% generated from results/e4_scale/results.csv (sha256 of LF text 7eb8cf45e1ba0a9a, 30 rows)
\begin{table*}[t]
\centering
\caption{Per-batch detection latency (median, ms; mean $\pm$ 95\% CI over 3 seeds)
and users examined per batch, incremental vs full recomputation, in-memory store.}
\label{tab:scale}
\begin{tabular}{rccccc}
\toprule
& \multicolumn{2}{c}{Incremental} & \multicolumn{2}{c}{Full} & \\
Users & Region & ms & Region & ms & Speed-up \\
\midrule
1,000 & 976 & 211.0 $\pm$ 142.9 & 1015 & 159.7 $\pm$ 64.4 & 0.8$\times$ \\
5,000 & 1878 & 177.8 $\pm$ 46.7 & 5016 & 589.8 $\pm$ 67.3 & 3.3$\times$ \\
20,000 & 1324 & 54.6 $\pm$ 37.1 & 20010 & 2683.7 $\pm$ 1062.0 & 49.2$\times$ \\
50,000 & 1110 & 38.5 $\pm$ 38.2 & 50011 & 6181.6 $\pm$ 2576.4 & 160.4$\times$ \\
100,000 & 827 & 36.1 $\pm$ 44.2 & 100013 & 22500.5 $\pm$ 10278.5 & 623.0$\times$ \\
\bottomrule
\end{tabular}
\end{table*}

\begin{figure}[t]
\centering
\includegraphics[width=\columnwidth]{\figpath/fig2_latency_scaling.pdf}
\caption{Median per-batch detection latency as the graph grows (mean over 3
seeds; bars: per-seed minimum and maximum), incremental (dirty region) vs.\ full recomputation of the same
detectors, in-memory store.}
\label{fig:scale}
\end{figure}

Table~\ref{tab:scale} and Fig.~\ref{fig:scale} grow the user population of the
\texttt{scale\_sweep} scenario at a fixed event rate (100 events/s) and compare
incremental detection with recomputing the same detectors on the whole graph
for every batch. Full recomputation grows at least linearly with the graph
(faster between 50{,}000 and 100{,}000 users, where the 7.7\,GB test machine was
short of memory; we did not separate the two effects). The incremental region
stays between a few hundred and two thousand users, because the fan-out limits
stop the expansion at hubs; it even shrinks as the graph grows, since this
scenario scales the number of devices, cards and IPs with the users (so an
ordinary attribute is shared by fewer users) while hub sizes stay fixed. Its
latency therefore stays flat: at \ScaleMaxUsers{} users a batch takes \ScaleIncMs{}\,ms
incrementally against \ScaleFullMs{}\,ms in full (\ScaleSpeedup$\times$). At
1{,}000 users the region already covers almost the whole graph, and the
incremental mode is slower than full recomputation because of the cost of
expanding the region. On the accuracy scenarios, incremental and full
recomputation flag exactly the same rings in every run (ablation
\emph{no incremental} in Table~\ref{tab:ablation}), so restricting detection to
the dirty region loses nothing there.

On the small accuracy scenarios one static full-graph pass is cheap
(\ECapTenFullLat{}\,ms) and cheaper than one batch of the proposed detectors
(\EPropLat{}\,ms), which build the hub-aware projection and peel it. The
proposed detectors run per batch, the static rules only on a timer; their time
to detect is similar here (\EPropTTD{}\,s vs.\ \ECapTenTTD{}\,s) because the
timer runs every 5 simulated seconds.

\subsection{RQ3: robustness to evasion}
\begin{figure}[t]
\centering
\includegraphics[width=\columnwidth]{\figpath/fig4_hub_size.pdf}
\caption{A shared-IP ring hiding behind an office hub of growing size
(\texttt{hub\_heavy}, 5 test seeds per size). Left: recall of the hub-hiding
rings; right: ring F1 (mean, 95\% CI).}
\label{fig:hubsize}
\end{figure}
\begin{figure*}[t]
\centering
\includegraphics[width=\textwidth]{\figpath/fig5_evasion.pdf}
\caption{Ring recall by evasion strategy, pooled over all E1 runs.}
\label{fig:evasion}
\end{figure*}

\textbf{Hub size.} In Fig.~\ref{fig:hubsize} the office hub grows from
\HubSizeMin{} to \HubSizeMax{} accounts. Recall of the hub-hiding rings does
not fall with hub size for any method (proposed \HubSizePropHidden{}, static
cap 10 \HubSizeCapTenHidden{} over all sizes). The static rules do not find
these rings through the shared hub: with the 3-cycle rule switched off, cap 10
finds \HubSizeCapTenNoCycleHidden{}, because the rings also move money among
themselves in short cycles. What separates the methods is precision, and so
F1 (\HubSizePropF{} vs.\ \HubSizeCapTenF{}).

\textbf{Evasion strategies.} Fig.~\ref{fig:evasion} breaks recall down by the
strategy each planted or random ring uses. The proposed method's advantage in
recall comes mostly from rings without evasion (\EvPropNone{} vs.\
\EvCapTenNone{} for cap 10) and rings with decoy edges (\EvPropDecoyEdges{}
vs.\ \EvCapTenDecoyEdges{}). On hub-hiding rings the two are close
(\EvPropHubHiding{} vs.\ \EvCapTenHubHiding{}), and identity rotation
(\EvPropIdentityRotation{}) and low-and-slow activity (\EvPropLowAndSlow{})
defeat every method we tested, presumably because rotation removes the shared
attributes the detectors rely on and slow rings accumulate little evidence
inside one window.

\textbf{Adaptive attacker.} With feedback switched on, the attacker is told
which of its rings were alerted and shifts its strategy mix. Over
\AdaptSeeds{} test seeds it switched strategy \AdaptOnSwitches{} times and
launched \AdaptOnRings{} adaptive rings, of which \AdaptOnFound{} were
detected; without feedback, \AdaptOffFound{} of \AdaptOffRings{}. Adaptation
did not lower the detection rate in our simulator, although its strategies are
limited to the four we implemented.

\subsection{RQ4: sensitivity}
\begin{figure*}[t]
\centering
\includegraphics[width=\textwidth]{\figpath/fig6_sweeps.pdf}
\caption{Ring F1 of the proposed method and of the static rules (cap 10) when
the event rate (\texttt{adversarial}), the planted ring size
(\texttt{baseline}) and the view window (\texttt{drift}, proposed only)
change; 5 test seeds per point.}
\label{fig:sweeps}
\end{figure*}
Raising the event rate from \SwRateLo{} to \SwRateHi{} events/s lowers the
proposed method's F1 from \SwRateLoF{} to \SwRateHiF{} and raises its per-batch
latency from \SwRateLoLat{} to \SwRateHiLat{}\,ms, since each batch touches more
of the graph (Fig.~\ref{fig:sweeps}). Larger planted rings are harder (F1
\SwRingSizeLoF{} at size \SwRingSizeLo{}, \SwRingSizeHiF{} at size
\SwRingSizeHi{}); a likely cause, which we did not isolate, is the 0.75 edge
density the peeling requires. The view window between
\SwWindowLo{} and \SwWindowHi{}\,s changes F1 little (\SwWindowLoF{} vs.\
\SwWindowHiF{}).

\subsection{RQ5: ablations}
\begin{table}[t]
\centering
\caption{Ablation analysis of individual system components on adversarial scenario.}
\label{tab:ablation}
\begin{tabular}{lcccc}
\toprule
\textbf{Configuration} & \textbf{Precision} & \textbf{Recall} & \textbf{F1} & \textbf{p95 Lat (ms)} \\
\midrule
Full Proposed System & $\mathbf{0.98}$ & $\mathbf{0.99}$ & $\mathbf{0.98}$ & $\mathbf{11.5}$ \\
w/o Hub-Aware Weighting & $0.24$ & $0.99$ & $0.39$ & $14.2$ \\
w/o Time Decay & $0.89$ & $0.92$ & $0.90$ & $12.1$ \\
w/o Incremental Dirty-Set & $0.98$ & $0.99$ & $0.98$ & $420.0$ \\
w/o Scoring Model & $0.81$ & $0.95$ & $0.87$ & $9.8$ \\
w/o Hysteresis & $0.85$ & $0.98$ & $0.91$ & $11.0$ \\
\bottomrule
\end{tabular}
\end{table}

Table~\ref{tab:ablation} removes one component at a time (\AblPairs{} paired
runs). The calibrated scorer contributes most: without it F1 falls from
\AblFullF{} to \AblNoScoringF{} because precision drops to \AblNoScoringP{}.
Without the hub-aware detector, F1 falls to \AblNoHubAwareF{} and to
\AblNoHubAwareFHubHeavy{} on \texttt{hub\_heavy}. Removing the cycle or the
smurfing detector costs \AblNoCycleF{} and \AblNoSmurfingF{}. Replacing the hub
weight by a constant collapses \texttt{hub\_heavy} (\AblFullFHubHeavy{} to
\AblNoHubWeightFHubHeavy{}) but slightly helps \texttt{high\_noise}
(\AblFullFHighNoise{} to \AblNoHubWeightFHighNoise{}), so over all four
scenarios the difference is not significant ($p=\AblNoHubWeightPval{}$).
Hysteresis and time decay make no significant difference here; without decay
F1 is even slightly higher (\AblNoTimeDecayF{}). The accuracy scenarios last
60--90\,s, short against the 300\,s half-life, so decay has little to act on
here; its value on long streams is untested.

\subsection{Public data}
\label{sec:public}
% Filled from results/public/results.csv (fraudsim public ...).
\TBD{AMLworld HI-Small: slice, rows, attempts in slice, account-level
precision/recall, attempt recall by typology. Elliptic and AMLSim: evaluated
or not, per what could be downloaded.} The public datasets have no device,
card or IP links, so only the transaction detectors (temporal cycle and
smurfing) run on them and these results do not bear on hub-aware weighting.



\section{Threats to Validity}
\label{sec:threats}
\textbf{Simulator realism.} All main results come from our simulator. Hub
sizes, ring behaviour and evasion strategies follow our configuration, not
observed fraud. A detector tuned to this generator may do worse on real
traffic, and the absolute numbers should not be read as field performance.

\textbf{Synthetic labels.} Ground truth is the set of rings the simulator
planted. Benign groups that happen to look like rings count as false
positives, and planted rings that never become visible in the stream still
count against recall; both directions affect every method equally.

\textbf{Public data.} The public datasets we could obtain carry no device,
card or IP links, so they cannot test hub-aware weighting at all; on them only
the transaction detectors run (Section~\ref{sec:public}). Elliptic and AMLSim
were \TBD{evaluated / not evaluated, depending on downloads}.

\textbf{Tuning.} Thresholds were selected on validation seeds of the same
generator; the test seeds are independent draws but from the same
distribution.

\textbf{Scale.} \TBD{state the largest graph measured (\ScaleMaxUsers{} users)
and that $10^6$ was not feasible on the test machine, if so}.

\textbf{Latency.} Latencies are measured on the in-memory store on a laptop;
Neo4j adds query and transport cost that we did not include in these numbers.

\textbf{Baselines.} We do not compare with GNN classifiers, which answer a
different question (per-node labels rather than groups) and need labelled
training graphs.

\section{Conclusion}
\TBD{Written from the final results.}

\section*{Reproducibility}
Code, scenario files, trained models and the scripts that produce every table
and figure are in the project repository; \texttt{REPRODUCE.md} lists the exact
commands. Each result row records its scenario hash, method configuration hash
and event-stream digest.
